% Bagean: sajarah
% Bab: biografi
% Pérangan: georg-cantor

\documentclass[../../../include/open-logic-section]{subfiles}

\begin{document}

\olfileid{his}{bio}{can}

\olsection{Georg Cantor}

\olphoto{cantor-georg}{Georg Cantor}

Sawijining biografi wiwitan ngenani Georg Cantor (\textsc{gay}-org
\textsc{kahn}-tor) tau ngandharaké yèn dhèwèké lair lan ditemokaké
ing kapal kang lagi lelayaran menyang Saint Petersburg, Rusia,
lan wong tuwané ora dingertèni. Nanging crita iku ora bener,
senajan Cantor pancèn lair ing Saint Petersburg ing taun 1845.

Cantor nggayuh gelar doktor matématika ing Universitas Berlin
ing taun 1867. Dhèwèké misuwur amarga karyané ing téori
himpunan, lan dianggep minangka pangadeg téori himpunan
dadi dhisiplin panlitèn kang mandiri.
Dhèwèké wong kapisan kang mbuktèkaké yèn ana himpunan
tanpa wates kang ukurané béda-béda.
Téori-téoriné, mligi téoriné ngenani tanpa wates,
nyebabaké akèh debat ing kalangan matématikawan
wektu iku, lan karyané dadi pasulayan.

Kapercayan agama Cantor lan karya matématikané raket banget
sesambungané; dhèwèké malah ngaku yèn téori bilangan
transfinit diwènèhaké langsung marang dhèwèké déning
Gusti Allah. Ing mangsa pungkasan uripé, Cantor ngalami
gangguan kesehatan mental. Wiwit taun 1894, lan saya kerep
nalika nyedhaki pungkasan uripé, Cantor dirawat ing rumah
sakit. Kritik abot marang karyané, kalebu pedhoté
sesambungan karo matématikawan Leopold Kronecker,
nyebabaké dhèwèké depresi lan kelangan minat marang
matématika. Nalika ngalami depresi, Cantor ngalihaké
kawigatèné menyang filsafat lan sastra, lan malah
nerbitaké téori yèn Francis Bacon iku panulisé
lakon-lakon Shakespeare.

Cantor séda tanggal 6 Januari 1918 ing sanatorium
ing Halle.

\begin{reading}
Kanggo biografi Cantor kang luwih jangkep, delengen
\citet{Dauben1990} lan \citet{Grattan-Guinness1971}.
Pamanggih Cantor kang radikal uga diandharaké ing acara
BBC Radio 4 \emph{A Brief History of
  Mathematics} \citep{Sautoy2014}. Yèn panjenengan
kepéngin ngrungokaké téori-téori Cantor awujud rap,
delengen \citet{Rose2012}.
\end{reading}

\end{document}
